% Hafta 6 — Kanca (hook) nasıl algılanır?
% Derleme: py -3.12 tools/latex/derle.py docs/week-6/assets/h06-10-hook.tex
\documentclass[tikz,border=8pt]{standalone}
\usepackage{cen429-cizim}
\begin{document}
\begin{tikzpicture}
  \node[baslik] (b) at (0,0) {Kanca (hook) nasıl algılanır?};
  \node[altbaslik] at (b.south west) {kanca, fonksiyonun ilk baytlarını değiştirir};
  \node[iyikutu=48mm, anchor=north west] (n1) at (0,-2.4)
    {\kb{Normal fonksiyon}\\beklenen prolog baytları};
  \node[kotukutu=48mm, below=14mm of n1] (n2)
    {\kb{Kancalı fonksiyon}\\ilk baytlar $\rightarrow$ jmp trampoline};
  \node[kutu=48mm, anchor=west] (n3) at ($(n1.east)!0.5!(n2.east)+(24mm,0)$)
    {\kb{DENETİM}\\prolog baytlarını karşılaştır\\dlsym+dladdr ile\\fonksiyonun geldiği .so'ya bak};
  \draw[iyiok] (n1.east) -- (n3.west |- n1.east);
  \draw[kotuok] (n2.east) -- (n3.west |- n2.east);
  \node[fit=(n1)(n2)(n3), inner sep=0pt] (grp) {};
  \node[dip, text width=155mm, anchor=north west] at ($(grp.south west)+(0,-8mm)$)
    {libc/vDSO dışında bir modülden geliyorsa kanca vardır (LD\_PRELOAD).};
\end{tikzpicture}
\end{document}
